\section{Conclusión}

En este trabajo hemos construido una teoría capaz de describir formalmente el Go clásico a partir de nociones elementales \textit{dónde, quiénes, qué y cómo}, sustentada en la teoría de conjuntos. A lo largo de esta construcción, las reglas del juego emergen de manera natural y conceptos intuitivos adquieren una expresión precisa y manipulable.

Sin embargo, el valor principal de este marco no se limita a formalizar el Go. Como se exploró en la Parte~\ref{parte3}, la estructura desarrollada es suficientemente general como para abarcar una familia más amplia de juegos basados en la idea de rodear y capturar. Esto sugiere la existencia de un objeto matemático más abstracto, al que denominamos \textsc{Juegos de Rodear} (\textsc{JR}). Del cual el Go es solo un caso particular.

Quedan abiertos varios caminos: identificar qué propiedades son exclusivas del Go y cuáles pertenecen a toda la familia, formalizar esta familia y ver hasta donde podemos llevarla, construir el Go infinito o bien, intentar conectar esta teoría con la \textit{teoría de juegos combinatorios}.

Si estas páginas han logrado transmitir aunque sea una fracción de la belleza que se esconde detrás de un tablero y unas cuantas piedras, entonces habrán cumplido su propósito. Además, si usted ya juega al Go, o si la lectura de estas páginas le ha motivado a intentarlo (un milagro), le dejo mi perfil de \href{https://online-go.com/user/view/1220808}{OGS}, una de las tantas plataformas para jugar al Go por internet. En caso de pertenecer al primer grupo, no espere de mí un nivel magistral; me defiendo, pero hasta ahí. En caso de ser del segundo grupo, cree una cuenta e invíteme a una partida, y gustosamente intentaré enseñarle los fundamentos (al menos tal como los entiendo) y señalarle qué recursos existen — véase la \href{https://www.go.org.ar/}{Asociación Argentina de Go}, por ejemplo — para dar sus primeros pasos.